\documentclass[11pt]{article}
\usepackage[utf8]{inputenc}
\usepackage{amsmath}
\usepackage{amsfonts}
\usepackage{amssymb}
\usepackage{bm}
\usepackage{geometry}
\usepackage{xcolor}
\usepackage{booktabs}
\usepackage{algorithm}
\usepackage{algorithmic}


\geometry{margin=0.8in}

\title{Synthetic Neuronal Network Topology Generation:\\
Distance-Dependent Sparse Connectivity with Dale's Law}
\author{Jan Balewski, NERSC}
\date{\today}

\begin{document}

\maketitle

\begin{abstract}
We describe the pipeline implemented in \texttt{gen\_daleMatrices4.py}:
(1)~build a spatially embedded, Dale-law-compliant sparse network
$A_{\text{true}}$ with target spectral radius $R<1$ and a bias vector $B$;
(2)~simulate spike counts under a Poisson GLM, either
\textbf{Model~A} (lag-1, no off-diagonal memory) or
\textbf{Model~B} (lag-$M$: diagonal lag-1 drive plus an off-diagonal
history filtered by a damped-oscillator kernel $\kappa_\ell$).
A third variant (Model~C: additional diagonal adaptation kernel) is
discussed separately but is \emph{not} implemented in this script.
\end{abstract}

\tableofcontents
\newpage

% ===============================================================
\section{Overview of the Method}
\label{sec:overview}
% ===============================================================

The objective is to produce synthetic spike-count time series
$Y_t\in\mathbb{Z}_{\geq 0}^{N}$ ($t=1,\dots,T$) with configurable
spatial connectivity and Poisson GLM dynamics.

\paragraph{Stage 1 --- Network construction (Section~\ref{sec:generation_AB}).}
$N$ neurons are sampled from a regular grid on $[0,L]\times[0,H]$ and
connected by distance-dependent out-degree sampling with Dale's law.
Outputs: $A_{\text{true}}$ with $\rho(A_{\text{true}})=R$ and bias $B$.

\paragraph{Stage 2 --- Spike generation (Section~\ref{sec:spike_gen}).}
\textbf{Model~A:} lag-1 recurrence on all couplings.
\textbf{Model~B:} diagonal term uses $Y_{t-1}$ only; off-diagonal
block uses lags $1,\dots,M$ with kernel $\kappa_\ell$ (Section~\ref{sec:lagM_single}).
Model~C (Section~\ref{sec:lagM_two}) is not implemented in code.

\paragraph{Notation conventions.}
$A_{\text{true}}$ denotes the connectivity matrix throughout.\\
$A_{\text{true}}^{\mathrm{diag}}=\operatorname{diag}
((A_{\text{true}})_{11},\dots,(A_{\text{true}})_{NN})$ is its diagonal
and $A_{\text{true}}^{\mathrm{off}}=A_{\text{true}}-
A_{\text{true}}^{\mathrm{diag}}$ its off-diagonal part.
In Model~B (Section~\ref{sec:lagM_single}), $A^{\mathrm{diag}}\in\mathbb{R}^N$
denotes the \emph{vector} of diagonal entries $(A_{\text{true}})_{ii}$,
parallel to the matrix shorthand $A^{\mathrm{off}}$.
$Y_t\in\mathbb{Z}_{\geq 0}^{N}$ is the vector of spike counts in
time bin $t$, where each bin has width $\Delta t$ seconds.
$B\in\mathbb{R}^{N}$ is the bias (log-baseline-rate) vector.
$\eta_{\mathrm{clip}}=5$ is a saturation threshold that prevents
numerical overflow in the exponential link function.
Column convention: $({A_{\text{true}}})_{ij}$ is the weight from
presynaptic neuron $j$ to postsynaptic neuron $i$.


% ===============================================================
\section{Generation of $A_{\text{true}}$ and $B$}
\label{sec:generation_AB}
% ===============================================================

\subsection{Parameters}

The script takes (among others): $N$ and $N_{\text{exc}}$ (default
$0.8\,N$); rectangular domain $[0,L]\times[0,H]$; grid spacing
$d_{\text{min}}>0$ (\texttt{--placement\_min\_dist}); distance-kernel
exponent $\delta_{\mathrm{ker}}\in\{1,2\}$ (third entry of
\texttt{--placement\_H\_L\_delta}); out-degree range from two fractions
$p_{\text{lo}}<p_{\text{hi}}$ via
$k_{\text{min}}=\max(1,\lfloor p_{\text{lo}}N\rfloor)$,
$k_{\text{max}}=\min(N-1,\lfloor p_{\text{hi}}N\rfloor)$ (defaults
$0.05$, $0.2$); weight spread $v\in(0,1)$; target spectral radius
$R\in(0,1)$.

\medskip\noindent
\textbf{Variables.}
$\mathbf{P}\in\mathbb{R}^{N\times 2}$: neuron spatial positions, row
$i=(x_i,y_i)$, sorted by $(x_i,y_i)$ lexicographically after placement.
$\mathbf{D}\in\mathbb{R}^{N\times N}$: pairwise Euclidean distances.
$\mathbf{V}\in\mathbb{R}^{N\times N}$: distance-kernel affinity matrix.
$\mathbf{E}\in\{-1,0,+1\}^{N\times N}$: signed topology matrix;
$E_{ij}\neq 0$ means neuron $j$ drives neuron $i$.
$\tau_j\in\{E,I\}$: type of neuron $j$;
$\sigma_j=+1$ (exc) or $-1$ (inh).
$k_j$: out-degree of neuron $j$.
$r=N_{\text{exc}}/N_{\text{inh}}=4$: E-I ratio.

\subsection{Algorithm}

\begin{algorithm}[H]
\caption{Synthetic Neuronal Network Topology}
\begin{algorithmic}[1]
\STATE \textbf{Phase 1:} Construct placement grid; sample $N$ nodes;
       assign types E/I and signs $\sigma_j$
\STATE \textbf{Phase 2:} Compute distance matrix $\mathbf{D}$ and affinity
       matrix $\mathbf{V}$
\STATE \textbf{Phase 3:} Sample sparse signed topology $\mathbf{E}$
\STATE \textbf{Phase 4:} Assign edge weights; enforce E-I ratio and
       rescale to spectral radius $R$
\STATE \textbf{Phase 5:} Generate bias vector $B$
\end{algorithmic}
\end{algorithm}

\subsubsection{Phase 1: Grid-Based Neuron Placement and Type
               Assignment}

\paragraph{Step 1.1 --- Construct placement grid.}
Tile $[0,L]\times[0,H]$ with spacing $d_{\text{min}}$:
\begin{equation}
  \mathcal{G} =
  \bigl\{(m\,d_{\text{min}},\;n\,d_{\text{min}}):\;
  m=0,\dots,\lfloor L/d_{\text{min}}\rfloor,\quad
  n=0,\dots,\lfloor H/d_{\text{min}}\rfloor\bigr\}.
\end{equation}
\begin{equation}
  N_{\mathcal{G}} =
  \bigl(\lfloor L/d_{\text{min}}\rfloor+1\bigr)
  \bigl(\lfloor H/d_{\text{min}}\rfloor+1\bigr).
\end{equation}
A valid configuration requires $N_{\mathcal{G}}\geq N$; if not
satisfied, $d_{\text{min}}$ must be reduced or $L$ increased.

\paragraph{Step 1.2 --- Sample neuron positions.}
Draw $N$ nodes uniformly at random \emph{without replacement} from
$\mathcal{G}$ and discard the remaining $N_{\mathcal{G}}-N$ nodes.
Store accepted positions as rows of $\mathbf{P}$, sorted by
$(x,y)$ (lexicographic).

\paragraph{Step 1.3 --- Assign neuron types.}
Sample $N_{\text{exc}}$ neuron indices without replacement as
excitatory; the remainder are inhibitory. Define:
\begin{equation}
  \sigma_j =
  \begin{cases} +1 & \tau_j=E,\\ -1 & \tau_j=I.\end{cases}
\end{equation}

\subsubsection{Phase 2: Distance $D$ and Affinity Matrix $V$}

Pairwise distances use Euclidean geometry; for the affinity, the
code sets $D_{ii}=\infty$ on the diagonal so self-entries are excluded
before
\begin{equation}
  V_{ij}=D_{ij}^{-\delta_{\mathrm{ker}}},\quad i\neq j;\qquad
  V_{ii}=0.
  \label{eq:dist_kernel}
\end{equation}
The saved distance matrix in \texttt{.simTruth.npz} uses zeros on the
diagonal.
$\delta_{\mathrm{ker}}\in\{1,2\}$ controls decay steepness; sampled
grid nodes satisfy spacing $\ge d_{\text{min}}$ so pairwise distances
among distinct neurons need not hit~$0$.

\subsubsection{Phase 3: Sparse Topology Sampling}

For each presynaptic neuron $j=1,\dots,N$ independently:
\begin{enumerate}
  \item Sample out-degree
        $k_j\sim\mathrm{Uniform\_int}(k_{\text{min}},k_{\text{max}})$.
  \item Normalize affinity over valid targets:
  \begin{equation}
    p_{ij}=\frac{V_{ij}}{\sum_{i'\neq j}V_{i'j}},\qquad i\neq j.
  \end{equation}
  \item Draw $k_j$ distinct postsynaptic targets
        $\mathcal{T}_j\subset\{1,\dots,N\}\setminus\{j\}$ without
        replacement according to $\{p_{ij}\}$.
\end{enumerate}
Build the signed topology matrix $E$ with negative self-loops:
\begin{equation}
  E_{ij}=
  \begin{cases}
    \sigma_j & i\in\mathcal{T}_j,\;i\neq j,\\
    -1       & i=j,\\
    0        & \text{otherwise.}
  \end{cases}
  \label{eq:topology}
\end{equation}

\subsubsection{Phase 4: Weight Assignment and Spectral Normalization}

\paragraph{Raw weights.}
Draw unsigned weights magnitudes
$U_{ij}\sim\mathrm{Uniform}(1-v,1+v)$ independently for each
non-zero entry of $\mathbf{E}$, then scale by neuron type:
\begin{equation}
  W_{ij}=
  \begin{cases}
    +U_{ij}    & E_{ij}=+1,\\
    -r\,U_{ij} & E_{ij}=-1,\;i\neq j,\\
    -U_{ii}    & i=j,\\
    0           & E_{ij}=0.
  \end{cases}
  \label{eq:weights}
\end{equation}

\paragraph{Spectral normalization.}
A single global rescaling sets $\rho(A_{\text{true}})$ to exactly $R$:
\begin{equation}
  \rho_0=\max_k|\lambda_k(W)|,\qquad
  A_{\text{true}}\leftarrow\frac{R}{\rho_0}\,W,
  \label{eq:spectral}
\end{equation}
where $\lambda_k(W)$ are the (possibly complex) eigenvalues of $W$.
Assume $\rho_0>0$.
Rescaling preserves all signs and the zero pattern, so Dale's law,
E-I ratio, and negative self-loops are all maintained.
For $R<1$, contractivity of $A_{\text{true}}$ is a sufficient
condition for stationarity of the Poisson GLM.

\subsubsection{Phase 5: Bias Vector Generation}

Let $f_{\mathrm{idle}}\in[f_{\min},f_{\max}]$ be a \emph{target idle
firing-rate band} in Hz (argument \texttt{idleRate}, default
$[15,30]$).  Independent draws
\begin{equation}
  B_i \sim \mathrm{Uniform}\bigl(\ln(R f_{\min}),\,\ln(R f_{\max})\bigr),
  \label{eq:bias}
\end{equation}
then all excitatory indices receive a fixed downward shift
$B_i\leftarrow B_i-2$ (hard-coded in the script).  Inhibitory
components are unchanged.  This matches \texttt{set\_flat\_selfSpiking}
in code.

\subsection{Output Summary}

\begin{table}[h]
\centering
\begin{tabular}{llll}
\toprule
Array & Shape & Type & Description \\
\midrule
$\mathbf{P}$          & $N\times 2$ & float
  & Neuron positions $(x_i,y_i)$, sorted by $(x,y)$ \\
$\mathbf{D}$          & $N\times N$ & float
  & Pairwise distances; diagonal $0$ in output \\
$\tau$                & $N$         & int
  & Type labels: $0$=exc, $1$=inh \\
$\sigma$              & $N$         & $\pm1$
  & Synaptic signs \\
$\mathbf{E}$          & $N\times N$ & $\{-1,0,+1\}$
  & Signed topology (Eq.~\ref{eq:topology}) \\
$k$                   & $N$         & int
  & Out-degree, $k_j\in[k_{\text{min}},k_{\text{max}}]$ \\
$A_{\text{true}}$     & $N\times N$ & float
  & Weighted connectivity, $\rho(A_{\text{true}})=R$ \\
$B$                   & $N$         & float
  & Bias vector (Eq.~\ref{eq:bias}) \\
\bottomrule
\end{tabular}
\caption{Output arrays of Stage~1.
$A_{\text{true}}$ and $B$ are passed to the spike generator
(Section~\ref{sec:spike_gen}).}
\label{tab:output}
\end{table}

% ===============================================================
\section{Generation of Spike Trains}
\label{sec:spike_gen}
% ===============================================================

Given $A_{\text{true}}$, $B$, time-bin width $\Delta t$, horizon
$T$, and $\eta_{\mathrm{clip}}=5$ (hard-coded), Models~A and~B are
implemented.  Bin $t$ indexes rows $Y_t$ (code uses $t=0,\dots,T-1$).
\textbf{Model~A:} $Y_0\sim\mathrm{Poisson}(\exp(B)\Delta t)$, then
Eq.~\eqref{eq:lag1} for $t\ge 1$.
\textbf{Model~B:} the first $M=\texttt{mem\_lag\_steps}$ rows are set
to zero (pre-history); dynamics follow Section~\ref{sec:lagM_single}
for $t\ge 1$.

% ---------------------------------------------------------------
\subsection{Model A: Lag-1, No Memory, Not Bursting}
\label{sec:lag1}
% ---------------------------------------------------------------

The simplest model conditions only on the immediately preceding bin:
\begin{equation}
  \eta_t = A_{\text{true}}\,Y_{t-1}+B,
  \qquad
  \tilde{\eta}_t = \min(\eta_t,\;\eta_{\mathrm{clip}}),
  \qquad
  \mu_t = \exp(\tilde{\eta}_t)\,\Delta t,
  \qquad
  Y_t \sim \mathrm{Poisson}(\mu_t).
  \label{eq:lag1}
\end{equation}
Here $\min(\cdot)$ and $\exp(\cdot)$ act element-wise. 
Because $\rho(A_{\text{true}})<1$, the process admits a stationary
distribution.

% ---------------------------------------------------------------
\subsection{Model B: Lag-$M$, Single Off-Diagonal Kernel, Bursting}
\label{sec:lagM_single}
% ---------------------------------------------------------------

\subsubsection{Off-diagonal kernel $\kappa_\ell$ (damped oscillator)}

Let $\tau_{\mathrm{mem}}>0$ be the oscillation period in \textbf{seconds}
(\texttt{time\_kernel\_a\_q\_tau}; if $\tau_{\mathrm{mem}}\le 0$, the
code sets $\tau_{\mathrm{mem}}=\tfrac{4}{3}M\Delta t$).
Superscript ${}^{(s)}$ marks quantities expressed in \textbf{discrete steps}
(bins), as distinct from the same quantity measured in seconds.
Define $\tau^{(s)}=\mathrm{round}(\tau_{\mathrm{mem}}/\Delta t)\ge 1$, the oscillation period in bins.
Let $Q>0$ and $a_{\mathrm{mem}}>0$ (quality factor and amplitude).
Set $\omega=2\pi/\tau^{(s)}$ and $\gamma=\pi/(Q\,\tau^{(s)})$.
For lags $\ell=1,\dots,M$ with $M=\texttt{mem\_lag\_steps}$,
\begin{equation}
  \kappa_\ell =
  \begin{cases}
    a_{\mathrm{mem}}\,
    \exp(-\gamma \ell)\,\cos(\omega \ell),
    & 4\ell \le 3\tau^{(s)},\\[0.35em]
    0, & \text{otherwise.}
  \end{cases}
  \label{eq:kappa_offdiag_gen4}
\end{equation}
There is \emph{no} $\ell=0$ term and \emph{no} normalization of
$\kappa$ to unit $\ell^1$ mass; amplitude is set solely by
$a_{\mathrm{mem}}$.

\subsubsection{Evolution}

Write $A^{\mathrm{diag}}\in\mathbb{R}^N$ with
$(A^{\mathrm{diag}})_i=(A_{\text{true}})_{ii}$, and let
$A^{\mathrm{off}}=A_{\text{true}}$ with zero diagonal.
The off-diagonal spike-history vector is
$S_t=\sum_{\ell=1}^{M}\kappa_\ell\,Y_{t-\ell}$ (vector sum).
For $t\ge 1$,
\begin{equation}
  \eta_t
  = A^{\mathrm{diag}}\odot Y_{t-1} + A^{\mathrm{off}} S_t + B,
  \label{eq:modelB_drive}
\end{equation}
then $\tilde\eta_t=\min(\eta_t,\eta_{\mathrm{clip}})$,
$\mu_t=\exp(\tilde\eta_t)\,\Delta t$, and
$Y_t\sim\mathrm{Poisson}(\mu_t)$ element-wise as in
Eq.~\eqref{eq:lag1}.
The first $M$ rows $Y_0,\dots,Y_{M-1}$ are held at~$0$ (pre-history).

\paragraph{Relation to Model~A.}
Model~A is a separate code path (\texttt{--spike\_model A}); it is
\textbf{not} recovered by formally setting $M=1$ in Model~B.

\subsubsection{Stability}

\texttt{gen\_daleMatrices4.py} does not impose an extra analytic
stability test for Model~B beyond rescaling $A_{\text{true}}$ to
spectral radius $R<1$ (Eq.~\ref{eq:spectral}).  A sufficient
condition for a two-kernel extension appears in
Section~\ref{sec:lagM_two} (Model~C).


% ---------------------------------------------------------------
\subsection{Model C: Lag-$M$, Two Kernels, Bursting - not implemented}
\label{sec:lagM_two}
% ---------------------------------------------------------------

\subsubsection{Adaptation kernel for the diagonal}

Spike-frequency adaptation is modeled as a purely inhibitory,
exponentially decaying self-history kernel with time constant
$\tau_{\mathrm{adapt}}$ (in time bins) and amplitude
$\alpha_{\mathrm{adapt}}>0$.
Let $M_{\mathrm{d}}$ denote the number of lags retained for the
diagonal kernel (which may differ from $M$).
The unnormalized weights are:
\begin{equation}
  \tilde{h}^{\mathrm{diag}}_\ell 
  = \exp\!\left(-\frac{\ell}{\tau_{\mathrm{adapt}}}\right),
  \qquad \ell = 1,\dots,M_{\mathrm{d}}.
  \label{eq:adapt_kernel_raw}
\end{equation}
Normalizing and attaching the inhibitory sign:
\begin{equation}
  h^{\mathrm{diag}}_\ell 
  = -\,\alpha_{\mathrm{adapt}}\;
    \frac{\tilde{h}^{\mathrm{diag}}_\ell} 
         {\sum_{k=1}^{M_{\mathrm{d}}}
          \tilde{h}^{\mathrm{diag}}_k}, 
  \qquad \ell = 1,\dots,M_{\mathrm{d}}.
  \label{eq:adapt_kernel}
\end{equation}
All weights are negative, ensuring that a neuron's own past
activity is always suppressive.
The dimensionless amplitude $\alpha_{\mathrm{adapt}}$ controls the
overall strength of adaptation relative to the diagonal entries of
$A_{\text{true}}$.

\subsubsection{Evolution equations}

The history drive now has two kernel-filtered components:
\begin{equation}
  H_t
  = A_{\text{true}}^{\mathrm{diag}}
    \!\left(\sum_{\ell=1}^{M_{\mathrm{d}}}
    h^{\mathrm{diag}}_\ell\,Y_{t-\ell}\right) 
  + A_{\text{true}}^{\mathrm{off}}
    \!\left(\sum_{\ell=1}^{M}
    h^{\mathrm{off}}_\ell\,Y_{t-\ell}\right). 
  \label{eq:histC}
\end{equation}
Note that, unlike Model~B, the diagonal term now integrates over
$M_{\mathrm{d}}$ lags rather than using only $Y_{t-1}$.
Spike generation proceeds identically:
\begin{equation}
  \eta_t = H_t + B,
  \qquad
  \tilde{\eta}_t = \min(\eta_t,\;\eta_{\mathrm{clip}}),
  \qquad
  \mu_t = \exp(\tilde{\eta}_t)\,\Delta t,
  \qquad
  Y_t \sim \mathrm{Poisson}(\mu_t).
  \label{eq:lagM_gen_C}
\end{equation}

\subsubsection{Reduction to Model B}

Setting $M_{\mathrm{d}}=1$ and $\alpha_{\mathrm{adapt}}=1$ gives
$h^{\mathrm{diag}}_1=-1$, so the diagonal term becomes 
$-A_{\text{true}}^{\mathrm{diag}}Y_{t-1}$.
Since $(A_{\text{true}})_{ii}<0$ by construction
(negative self-loops), the product
$-A_{\text{true}}^{\mathrm{diag}}Y_{t-1}
 =|A_{\text{true}}^{\mathrm{diag}}|Y_{t-1}$
has the wrong sign.
To recover Model~B exactly, one must instead set
$\alpha_{\mathrm{adapt}}=-1$ (i.e.\ disable the adaptation sign
flip) or, equivalently, simply bypass the diagonal kernel and use
$A_{\text{true}}^{\mathrm{diag}}Y_{t-1}$ directly.
In practice, Model~C is not intended to reduce to Model~B; rather,
Model~B is the appropriate choice when adaptation is not modeled.

\subsubsection{Stability}

A sufficient condition for stability of Model~C is:
\begin{equation}
  \rho\!\bigl(A_{\text{true}}^{\mathrm{off}}\bigr)
  \cdot\sum_{\ell=1}^{M}|h^{\mathrm{off}}_\ell| 
  \;+\;
  \max_i\bigl|(A_{\text{true}})_{ii}\bigr|
  \cdot\sum_{\ell=1}^{M_{\mathrm{d}}}
        |h^{\mathrm{diag}}_\ell| 
  \;<\; 1.
  \label{eq:stabilityC}
\end{equation}
The first term equals
$\rho(A_{\text{true}}^{\mathrm{off}})$ by the off-diagonal kernel
normalization.
The second term equals
$\alpha_{\mathrm{adapt}}\max_i|(A_{\text{true}})_{ii}|$ by the
diagonal kernel normalization.
Both contributions must be kept small enough that their sum remains
below unity.

% ---------------------------------------------------------------
\subsection{Parameter Choices}
\label{sec:params}
% ---------------------------------------------------------------

The script defaults include $\Delta t=0.01$\,s and
\texttt{time\_kernel\_a\_q\_tau} $=(a_{\mathrm{mem}},Q,\tau_{\mathrm{mem}})=(1,3,0)$
(the third value $0$ triggers
$\tau_{\mathrm{mem}}=\tfrac{4}{3}M\Delta t$ in seconds when
\texttt{mem\_lag\_steps} $=M$ is set).
Choose $M$ large enough that non-zero $\kappa_\ell$ out to
$\ell\le \lfloor 3\tau^{(s)}/4\rfloor$ are retained
(Eq.~\ref{eq:kappa_offdiag_gen4}); the code prints the minimum
$M$ needed for full support.

Model~C-specific parameter guidance remains in
Section~\ref{sec:lagM_two}.

\newpage
% ---------------------------------------------------------------
\subsection{Notation Summary}
\label{sec:notation}
% ---------------------------------------------------------------

For reference, all symbols used in the spike generation models are
collected here.

\begin{table}[h]
\centering
\begingroup
\renewcommand{\arraystretch}{1.22}
\begin{tabular}{lll}
\toprule
Symbol & Domain & Description \\
\midrule
$N$    & $\mathbb{Z}_{>0}$ & Number of neurons \\
$T$    & $\mathbb{Z}_{>0}$ & Number of time bins \\
$\Delta t$ & $\mathbb{R}_{>0}$ &
  Time-bin width (seconds); default $0.01$\,s \\
$Y_t$  & $\mathbb{Z}_{\geq 0}^{N}$ &
  Spike-count vector in bin $t$ \\
$A_{\text{true}}$ & $\mathbb{R}^{N\times N}$ &
  Connectivity matrix from Section~\ref{sec:generation_AB} \\
$A_{\text{true}}^{\mathrm{diag}}$ & $\mathbb{R}^{N\times N}$ &
  Diagonal of $A_{\text{true}}$ (self-coupling) \\
$A_{\text{true}}^{\mathrm{off}}$ & $\mathbb{R}^{N\times N}$ &
  Off-diagonal of $A_{\text{true}}$ (cross-neuron coupling) \\
$A^{\mathrm{diag}}$ & $\mathbb{R}^{N}$ &
  Vector $(A_{\text{true}})_{ii}$; Model~B shorthand (Eq.~\ref{eq:modelB_drive}) \\
$B$    & $\mathbb{R}^{N}$ &
  Bias vector (log-baseline rate) \\
$\eta_{\mathrm{clip}}$ & $\mathbb{R}_{>0}$ &
  Clipping threshold; default $5$ \\
\midrule
\multicolumn{3}{l}{\textit{Off-diagonal kernel (Model~B; Model~C off-part)}} \\
$\tau_{\mathrm{mem}}$ & $\mathbb{R}_{>0}$ &
  Oscillation period (\textbf{seconds}); $\tau^{(s)}=\mathrm{round}(\tau_{\mathrm{mem}}/\Delta t)$ \\
$Q$    & $\mathbb{R}_{>0}$ &
  Quality factor (\texttt{Q\_mem}) \\
$\omega$ & $\mathbb{R}_{>0}$ &
  $2\pi/\tau^{(s)}$ \\
$\gamma$ & $\mathbb{R}_{>0}$ &
  $\pi/(Q\,\tau^{(s)})$ \\
$M$    & $\mathbb{Z}_{>0}$ &
  \texttt{mem\_lag\_steps}; history lags $\ell=1,\dots,M$ \\
$a_{\mathrm{mem}}$ & $\mathbb{R}_{>0}$ &
  Kernel amplitude (\texttt{amp\_mem}) \\
$\kappa_\ell$ & $\mathbb{R}$ &
  Eq.~\ref{eq:kappa_offdiag_gen4} \\
\midrule
\multicolumn{3}{l}{\textit{Diagonal adaptation kernel
  (Model~C only)}} \\
$\tau_{\mathrm{adapt}}$ & $\mathbb{R}_{>0}$ &
  Adaptation time constant (time bins); default $10$ \\
$\alpha_{\mathrm{adapt}}$ & $\mathbb{R}_{>0}$ &
  Adaptation amplitude; default $0.5$ \\
$M_{\mathrm{d}}$ & $\mathbb{Z}_{>0}$ &
  Number of diagonal history lags \\
$h^{\mathrm{diag}}_\ell$ & $\mathbb{R}_{\leq 0}$ & 
  Diagonal kernel weight at lag $\ell$
  (Eq.~\ref{eq:adapt_kernel}) \\
\bottomrule
\end{tabular}
\endgroup
\caption{Complete notation for the spike generation models
  (Section~\ref{sec:spike_gen}).
  All symbols from Section~\ref{sec:generation_AB}
  ($\mathbf{P}$, $\mathbf{E}$, $\sigma$, $\tau_j$, $r$, $R$,
  $\delta_{\mathrm{ker}}$, etc.)\ retain their definitions
  from that section.}
\label{tab:notation}
\end{table}

% = = = = = = = = = ==== 
% = = = = = = = = = ==== 
\appendix
\section*{\Huge Appendix}
\addcontentsline{toc}{section}{Appendix}
\vspace{1cm}

\subsubsection{Excitable memory kernel}

Experimental organoid recordings show brief synchronous bursts
separated by long, irregular quiescent intervals.  This pattern is
characteristic of an excitable system rather than an oscillator:
recent activity first promotes further firing (burst build-up)
and then strongly suppresses it (post-burst refractoriness).
We model this with a difference-of-exponentials kernel:
\begin{equation}
  h(\tau) = e^{-\tau/\tau_{\mathrm{exc}}}
          - \beta\,e^{-\tau/\tau_{\mathrm{inh}}},
  \qquad \tau>0,
\end{equation}
where $\tau_{\mathrm{exc}}\ll\tau_{\mathrm{inh}}$ and
$\beta\in(0,1)$.
The fast excitatory lobe ($\tau_{\mathrm{exc}}$, a few tens of
milliseconds) creates the burst; the slow inhibitory lobe
($\tau_{\mathrm{inh}}$, several hundred milliseconds) enforces a
prolonged refractory period.
The kernel crosses zero at
$\tau^{*}=\frac{\tau_{\mathrm{exc}}\,\tau_{\mathrm{inh}}}
               {\tau_{\mathrm{inh}}-\tau_{\mathrm{exc}}}
 \ln(1/\beta)$.
Because the inhibitory tail decays slowly, the network remains
quiescent until the suppression wears off and a noise fluctuation
triggers the next burst, producing irregular inter-burst intervals
without a fixed oscillation period.

Sampling at integer lags and normalizing gives the discrete
weights:
\begin{equation}
  \tilde{h}^{\mathrm{off}}_\ell
  = \exp(-\ell/\tau_{\mathrm{exc}})
  - \beta\,\exp(-\ell/\tau_{\mathrm{inh}}),
  \qquad
  h^{\mathrm{off}}_\ell
  = \frac{\tilde{h}^{\mathrm{off}}_\ell}
         {\sum_{k=1}^{M}|\tilde{h}^{\mathrm{off}}_k|},
  \qquad \ell=1,\dots,M.
  \label{eq:mem_kernel}
\end{equation}
A kernel support of $M\geq 3\tau_{\mathrm{inh}}$ suffices to
capture the exponential tail.


\subsubsection{Asymmetric damped-oscillator memory kernel}

Experimental organoid recordings show brief synchronous bursts
separated by long, irregular quiescent intervals.  The standard
damped-cosine kernel $h(\tau)=e^{-\gamma\tau}\cos(\omega\tau)$
produces symmetric positive and negative lobes, so the excitatory
and suppressive phases of each cycle have equal duration.  To
obtain short bursts followed by prolonged silences we introduce a
smooth, single-parameter frequency modulation that compresses the
positive (excitatory) lobe and stretches the negative
(suppressive) lobe.

Define the phase-modulated kernel
\begin{equation}
  h(\tau) = e^{-\gamma\tau}
            \cos\!\bigl(\omega\tau + \delta\sin(\omega\tau)\bigr),
  \qquad \tau>0,
  \label{eq:asym_kernel}
\end{equation}
with
\begin{equation}
  \omega = \frac{2\pi}{\tau_{\mathrm{mem}}},
  \qquad
  \gamma = \frac{\pi}{Q\,\tau_{\mathrm{mem}}},
  \label{eq:asym_params}
\end{equation}
exactly as in the symmetric case.  The new dimensionless parameter
$\delta\in[0,1)$ controls the asymmetry between the two lobes.

The instantaneous frequency of the argument is
$\dot\Phi(\tau)=\omega\bigl(1+\delta\cos(\omega\tau)\bigr)$.
Near the positive peak of the cosine
($\cos(\omega\tau)\approx+1$) the phase advances at rate
$\omega(1+\delta)$, compressing the positive lobe; near the
negative trough ($\cos(\omega\tau)\approx-1$) it advances at rate
$\omega(1-\delta)$, stretching the negative lobe.  The ratio of
time spent in the negative versus positive half-cycle is
approximately $(1+\delta)/(1-\delta)$.  Setting $\delta=0$
recovers the original symmetric kernel; $\delta=0.6$ makes the
negative lobe roughly four times longer than the positive one.

Sampling at integer lags and normalizing by the sum of absolute
values gives the discrete weights:
\begin{equation}
  \tilde{h}^{\mathrm{off}}_\ell
  = \exp(-\gamma\,\ell)\,
    \cos\!\bigl(\omega\,\ell+\delta\sin(\omega\,\ell)\bigr),
  \qquad
  h^{\mathrm{off}}_\ell
  = \frac{\tilde{h}^{\mathrm{off}}_\ell}
         {\sum_{k=1}^{M}|\tilde{h}^{\mathrm{off}}_k|},
  \qquad \ell=1,\dots,M.
  \label{eq:asym_mem_kernel}
\end{equation}
A kernel support of $M\geq 2\tau_{\mathrm{mem}}$ suffices to
capture the full tail.  Typical parameter choices are
$\tau_{\mathrm{mem}}=40$--$60$ bins, $Q=2$--$3$, and
$\delta=0.5$--$0.7$. 
\subsubsection{Evolution equations}

Decompose $A_{\text{true}}$ into diagonal and off-diagonal parts:
\begin{equation}
  A_{\text{true}}^{\mathrm{diag}}
  = \operatorname{diag}\!\bigl((A_{\text{true}})_{11},\dots,
    (A_{\text{true}})_{NN}\bigr),
  \qquad
  A_{\text{true}}^{\mathrm{off}}
  = A_{\text{true}} - A_{\text{true}}^{\mathrm{diag}}.
\end{equation}
The history drive depending on the past $M$ time bins is: 
\begin{equation}
  H_t
  = A_{\text{true}}^{\mathrm{diag}}\,Y_{t-1}
  + A_{\text{true}}^{\mathrm{off}}
    \!\left(\sum_{\ell=1}^{M}
    h^{\mathrm{off}}_\ell\,Y_{t-\ell}\right). 
  \label{eq:histB}
\end{equation}
The diagonal term uses only the most recent time bin (lag-1 self-coupling),
while the off-diagonal term integrates kernel-weighted spike history
over $M$ lags.
Spikes are then generated as in Model~A:
\begin{equation}
  \eta_t = H_t + B,
  \qquad
  \tilde{\eta}_t = \min(\eta_t,\;\eta_{\mathrm{clip}}),
  \qquad
  \mu_t = \exp(\tilde{\eta}_t)\,\Delta t,
  \qquad
  Y_t \sim \mathrm{Poisson}(\mu_t).
  \label{eq:lagM_gen}
\end{equation}

\end{document}
